\documentclass[UTF8,zihao=-4]{ctexart}

\usepackage[a4paper,margin=2.3cm]{geometry}
\usepackage{array}
\usepackage{booktabs}
\usepackage{longtable}
\usepackage{hyperref}
\usepackage{xcolor}
\usepackage{enumitem}
\usepackage{fancyhdr}
\usepackage{lastpage}

\hypersetup{
  colorlinks=true,
  linkcolor=black,
  urlcolor=blue
}

\setlength{\parindent}{2em}
\setlength{\parskip}{0.35em}
\renewcommand{\arraystretch}{1.25}
\setlist[itemize]{leftmargin=2.2em,itemsep=0.2em,topsep=0.2em}
\pagestyle{fancy}
\fancyhf{}
\fancyfoot[C]{第 \thepage 页 / 共 \pageref*{LastPage} 页}
\renewcommand{\headrulewidth}{0pt}
\renewcommand{\footrulewidth}{0pt}

\title{\heiti 托福阅读 Class 2 课堂笔记}
\author{}
\date{}

\begin{document}
\maketitle
\thispagestyle{fancy}

本节课继续围绕托福阅读中的词根、词缀和高频学术词展开。课堂重点包括表示“向外”的 \texttt{ex-}、表示“拉、拖、区域”的 \texttt{tract-}，以及与“空、糖、颜色、大小、双数”等语义相关的词根词缀。整理时保留课堂笔记中的核心词义，并对部分容易混淆的词作了归纳。

\section*{一、词根词缀与重点单词}

\texttt{antiqu-} 与“古代、古老”有关。\texttt{antiquated} 表示“陈旧的、过时的”，近义词可记为 \texttt{outmoded}。\texttt{in remote antiquity} 表示“在远古时期”。此外，\texttt{plunge} 表示“跳入、跌入、骤降”，\texttt{vicious cycle} 表示“恶性循环”。

\texttt{ex-} 常表示“向外、外部”。\texttt{extract} 或 \texttt{extracting} 表示“提取、抽取”。\texttt{exodus} 表示“大批离去、大规模迁徙”，\texttt{exosphere} 表示“外逸层”，\texttt{exotic} 表示“异国情调的”，\texttt{eccentric} 表示“离心的、古怪的”。\texttt{exogamy} 表示“外婚制、异族通婚”，对应形容词为 \texttt{exogamous}；\texttt{endogamy} 表示“内婚制、同族通婚”，对应形容词为 \texttt{endogamous}。

\texttt{tract-} 有“拉、拖、区域”等含义。\texttt{tract} 本身可以指“大片土地、地带”，也可指人体中的“道、束”，如 \texttt{the digestive tract} 表示“消化道”，\texttt{the respiratory tract} 表示“呼吸道”，\texttt{the vocal tract} 表示“声道”。相关单词包括：\texttt{subtract} 表示“减去”，\texttt{contract} 表示“收缩；合同、契约”，\texttt{detract} 表示“减损、贬低”，\texttt{distract} 表示“分散注意力”，\texttt{attract} 表示“吸引”，\texttt{retract} 表示“收回、撤回”，\texttt{protract} 表示“延长、拖延”，\texttt{abstract} 表示“抽象的；摘要”，\texttt{tractor} 表示“拖拉机”，\texttt{traction} 表示“牵引力；吸引力、关注度”。短语 \texttt{on the strength of} 可理解为 \texttt{on the basis of}，表示“基于、根据、凭借”。

与婚姻制度相关的词汇可以集中记忆。\texttt{monogamy} 表示“一夫一妻制”，\texttt{bigamy} 表示“重婚罪”，\texttt{polygamy} 表示“多配偶制，常指一夫多妻制”，\texttt{polyandrous} 表示“一雌多雄的、一妻多夫的”。\texttt{kin} 表示“亲属”，\texttt{kinship} 表示“亲属关系”，\texttt{be akin to} 表示“与某物相似”。\texttt{taboo} 表示“禁忌”。

\texttt{void / van / vac} 都可与“空、缺乏”有关。\texttt{be devoid of} 表示“全无、缺乏”，\texttt{void} 表示“空的、无效的；空缺”，\texttt{fill a void} 表示“填补空白”。\texttt{Vanity Fair} 可译为“名利场”，\texttt{evanescent} 表示“短暂易逝的”，\texttt{vacation} 表示“假期”，\texttt{vocation} 表示“职业、天职”，\texttt{vacuum} 表示“真空”，\texttt{vacuous} 表示“空洞的、无思想的”，如 \texttt{a vacuous romantic novel} 可译为“一本空洞的爱情小说”。\texttt{be susceptible to} 和 \texttt{be vulnerable to} 都可表示“易受……影响、伤害或攻击”。

\texttt{gluc- / glyc-} 与“糖”有关。\texttt{glycogenesis} 表示“糖原生成”，\texttt{glycated hemoglobin} 表示“糖化血红蛋白”，\texttt{hyperglycemia} 表示“高血糖”，\texttt{hypoglycemia} 表示“低血糖”，\texttt{glycerin} 表示“甘油”，\texttt{glucosamine} 表示“氨基葡萄糖”。

生物学阅读中还常见一组细胞结构词。\texttt{ribosome} 表示“核糖体”，\texttt{lysosome} 表示“溶酶体”，\texttt{proteasome} 表示“蛋白酶体”，\texttt{phagosome} 表示“吞噬体”，其中 \texttt{phag-} 与“吃、吞噬”有关；\texttt{endosome} 表示“内体、胞内体”。

\texttt{magn-} 与“大、伟大”有关。\texttt{magnificent} 表示“辉煌的、壮丽的”，\texttt{magnify} 表示“放大”，\texttt{magnanimous} 表示“宽宏大量的”，名词为 \texttt{magnanimity}；\texttt{magnitude} 表示“大小、量级、重要性”，\texttt{magniloquent} 表示“夸张的、夸夸其谈的”，\texttt{magnum opus} 表示“代表作、鸿篇巨制”。

\texttt{chrom-} 与“颜色”有关。\texttt{chromatic} 表示“彩色的”，\texttt{monochrome} 表示“单色的、黑白的”，\texttt{polychrome} 表示“多色的”，\texttt{chromosphere} 表示“色球层”。

\texttt{di-}、\texttt{dia-}、\texttt{du- / duo- / dual}、\texttt{bi-}、\texttt{twi-}、\texttt{ambi- / amphi-} 都可以帮助记忆与“二、双、两边”有关的词。\texttt{dioxide} 表示“二氧化物”，\texttt{dichotomy} 表示“二分法、对立”，\texttt{diphthong} 表示“双元音”，\texttt{diode} 表示“二极管”。\texttt{diagnose} 表示“诊断”，\texttt{dialogue} 表示“对话”，\texttt{diameter} 表示“直径”。\texttt{duet} 表示“二重奏”。\texttt{binary} 表示“二进制的、二元的”，\texttt{bilateral} 表示“双边的、双侧的”，\texttt{bicameral} 表示“两院制的”，\texttt{bilingual} 表示“双语的”，\texttt{biceps} 表示“二头肌”，\texttt{binoculars} 表示“双筒望远镜”。\texttt{twin} 表示“双胞胎”，\texttt{twilight} 表示“暮光、黄昏”。\texttt{ambivalent} 表示“矛盾的、摇摆不定的”，\texttt{ambidextrous} 表示“双手灵巧的”，\texttt{ambiguous} 表示“模棱两可的”，\texttt{amphibian} 表示“两栖动物”，\texttt{amphibious vehicle} 表示“水陆两用车”，\texttt{amphitheater} 表示“圆形剧场”，\texttt{amphipathic} 表示“两亲性的”。

\section*{二、易混词与补充表达}

\begin{longtable}{>{\raggedright\arraybackslash}p{4.8cm} >{\raggedright\arraybackslash}p{9.7cm}}
\toprule
\textbf{单词或表达} & \textbf{含义} \\
\midrule
\endfirsthead
\toprule
\textbf{单词或表达} & \textbf{含义} \\
\midrule
\endhead
\texttt{antiquated} & 陈旧的，过时的 \\
\texttt{outmoded} & 过时的，不流行的 \\
\texttt{in remote antiquity} & 在远古时期 \\
\texttt{plunge} & 跳入，跌入；骤降 \\
\texttt{vicious cycle} & 恶性循环 \\
\texttt{extract} & 提取，抽取；摘录 \\
\texttt{tract} & 大片土地；地带；人体的“道、束” \\
\texttt{subtract} & 减去 \\
\texttt{contract} & 收缩；合同，契约 \\
\texttt{detract} & 减损，贬低 \\
\texttt{distract} & 分散注意力 \\
\texttt{attract} & 吸引 \\
\texttt{retract} & 收回，撤回 \\
\texttt{protract} & 延长，拖延 \\
\texttt{abstract} & 抽象的；摘要 \\
\texttt{tractor} & 拖拉机 \\
\texttt{traction} & 牵引力；吸引力，关注度 \\
\texttt{on the strength of} & 基于，根据，凭借 \\
\texttt{derive from} & 源于，来自 \\
\texttt{stem from} & 源于，来自 \\
\texttt{originate from} & 起源于 \\
\texttt{arise from} & 由……产生 \\
\texttt{result from} & 由……导致，起因于 \\
\texttt{spring from} & 源于，突然来自 \\
\texttt{emanate from} & 发源于，散发自 \\
\texttt{taboo} & 禁忌 \\
\texttt{kin} & 亲属 \\
\texttt{kinship} & 亲属关系 \\
\texttt{be akin to} & 与……相似 \\
\texttt{be devoid of} & 全无，缺乏 \\
\texttt{be susceptible to} & 易受……影响、感染或损害 \\
\texttt{be vulnerable to} & 易受……攻击、伤害或影响 \\
\texttt{vacuous} & 空洞的，无思想的 \\
\texttt{vocation} & 职业，天职 \\
\texttt{vacation} & 假期 \\
\bottomrule
\end{longtable}

\end{document}
